% Fragment for the parent article. Requires amsmath, amssymb, amsthm.
% No documentclass or standalone bibliography is included.

\section{Several transverse harmonic marks and Gamma--Stieltjes gaps}
\label{gm:section}

The harmonic-gap continuation \cite{gm:gaps-source} leaves a precise
question after its single-marker theorem: can two independent transverse
derivatives be taken at adjacent deleted indices, while retaining the
prescribed relative continuation at nonintegral endpoints?  We answer its
first target,
\(\left.\partial_u\partial_v\mathcal H_{a,b}(s,u,v,t)\right|_{u=v=0}\),
and every analogous cluster of identical marks.  A finite set-partition
formula also determines the fully symmetric part of arbitrarily different
marks.  The partition coefficients are the classical symmetric-sum
coefficients of Hoffman \cite{gm:Hoffman1992,gm:Hoffman2018}.  The point
requiring proof here is their local realization by decorated relative
quotients, including products of Hurwitz primitives and their simultaneous
resonant derivatives.  An individual ordered word with unequal marks
retains information that the symmetric formula does not determine.

\subsection{Normalization and closure under finite decorations}

Throughout this section, \(\Re a,\Re b>0\), and all powers use the
analytic principal logarithm on the right half-plane.  For sufficiently
decaying lattice functions, define
\[
 Z_c[f_1,\ldots,f_d]
 =\sum_{n_1>\cdots>n_d\ge0}\prod_{j=1}^df_j(n_j+c),
 \qquad Z_c[\varnothing]=1.
\]
The relative quotient is fixed by
\begin{equation}\label{gm:relative}
 \mathcal H_{a,b}[w]
 =Z_b[w]-\sum_{j=1}^{|w|}Z_a[w_{\le j}]
                                  \mathcal H_{a,b}[w_{>j}],
 \qquad \mathcal H_{a,b}[\varnothing]=1.
\end{equation}
For power letters \(f_s(x)=x^{-s}\), write
\(\mathcal H_{a,b}(s_1,\ldots,s_d)\).  This is the relative
interpolant already used in the source; no new endpoint interpolation is
being chosen.

For clarity, its joint spectral entireness follows from the published
Ihara--Nakamura--Yamamoto theory \cite[Corollary~4.2 and
Proposition~5.3]{gm:INY}.  With deconcatenation convolution denoted by
\(*\), put \(\Phi_c(w)=\Psi(w;c-1)\) in their notation.  Their
formula is \(\Phi_c=Z_c^{-1}*Z_1\); consequently
\[
 \mathcal H_{a,b}=Z_a^{-1}*Z_b=\Phi_a*(\Phi_b)^{-1}.
\]
On each fixed word, convolution and inversion are finite polynomial
operations in the entire \(\Phi_c\)-values.  The harmonic product law
also follows from the character property
\cite[Theorem~5.2]{gm:INY}.  These identities initially hold on an open
absolute-convergence domain and then continue analytically.

Use the completed Hurwitz function
\begin{equation}\label{gm:Q}
 Q(u,x)=\zeta(u,x)-\frac1{u-1},\qquad
 Q(u,x+1)=Q(u,x)-x^{-u}.
\end{equation}
It is entire in \(u\) and holomorphic for \(\Re x>0\).
Write \(Q^{[r]}\) and \(\zeta^{[r]}\) for derivatives in the first
argument.  The Stieltjes convention is
\begin{equation}\label{gm:stieltjes}
 Q(1+\varepsilon,x)
  =\sum_{r\ge0}\frac{(-1)^r\gamma_r(x)}{r!}\varepsilon^r,
 \qquad Q^{[r]}(1,x)=(-1)^r\gamma_r(x).
\end{equation}

A decorated letter is \(x^{-s}\) multiplied by a polynomial in
finitely many \(Q(u,x)\) and \(Q(u,x+1)\).  Its \(Z_c\) and
\(\mathcal H_{a,b}\) are first defined by their actual convergent
series and \eqref{gm:relative}.  A common right half-plane in all
spectral variables exists.

\begin{lemma}[One tail with arbitrary other decorations]
\label{gm:one-tail}
Let \(w=(f_1,\ldots,f_d)\).  Denote by \(M_i(u)w\) the word with
\(f_i(x)\) replaced by \(f_i(x)Q(u,x)\), by \(I_j(u)w\) the
insertion of \(f_u\) after position \(j\), and by \(C_j(u)w\) the
replacement of \(f_j\) by \(f_jf_u\).  Then
\begin{equation}\label{gm:insertion}
 \mathcal H_{a,b}[M_i(u)w]
 =Q(u,a)\mathcal H_{a,b}[w]
  +\sum_{j=0}^{i-1}\mathcal H_{a,b}[I_j(u)w]
  +\sum_{j=1}^{i}\mathcal H_{a,b}[C_j(u)w].
\end{equation}
The other letters may already carry arbitrary finite decorations.
\end{lemma}

\begin{proof}
First replace \(Q\) by \(\zeta\) and work where all relevant sums
converge absolutely.  In the infinite decorated sum, the additional
Hurwitz-tail index is at least the index in position \(i\).  Its strict
positions are the \(i\) slots preceding that position; its equality
cases are the first \(i\) positions.  This gives the insertion and
merger terms with \(Z_c\) in place of \(\mathcal H_{a,b}\), without
an endpoint term.

Expand the relative recursion and group terms by their cut in the
original word.  At a cut before position \(i\), the insertions and
mergers contained in the prefix sum to
\(\zeta(u,a)Z_a[\text{prefix}]\), by the ordinary singleton stuffle
identity.  At a cut at or after position \(i\), the same counting
identity gives the decorated prefix.  The remaining suffix terms are
exactly the recursion for the right side of \eqref{gm:insertion}.
Triangular uniqueness in the word length proves that formula with
\(\zeta\).  Subtracting \(\mathcal H_{a,b}[w]/(u-1)\) gives the
stated completed formula.  This argument only uses pointwise products
of the other letters, so it applies unchanged when they are decorated.
\end{proof}

\begin{theorem}[Entire continuation of every finite decoration]
\label{gm:closure}
Every decorated relative word above continues jointly entirely in all
its spectral variables.  It has a terminating finite expansion into
bare relative words at affine sums of spectral arguments, with
coefficients polynomial in entire endpoint functions \(Q(U,a)\).
All fixed independent spectral derivatives are legitimate.
\end{theorem}

\begin{proof}
Use \eqref{gm:Q} to remove shifted factors.  Apply
Lemma~\ref{gm:one-tail} to one unshifted factor.  Every resulting
relative word contains one fewer \(Q\) factor: an insertion or merger
adds a bare power, while the remaining decorations are unchanged.
Induction ends with the entire bare relative words.  Each step agrees
with the initial convergent quotient, so the resulting entire expression
is its analytic continuation.  This proves both existence and
independence of the elimination order.  In the initial domain, Hurwitz
asymptotics uniformly on a slightly larger spectral compact set justify
termwise differentiation; elsewhere the derivatives follow from the
proved entire identities.  In particular the argument does not infer an
analytic identity from equality at integer endpoint differences.
\end{proof}

\subsection{The local symmetric-sum theorem}

For a finite separable polynomial \(P(y,z)\), define linearly
\begin{equation}\label{gm:E}
 \mathcal E_{a,b}^{s,t}[A(y)B(z)]
       =\mathcal H_{a,b}[f_sA,f_tB].
\end{equation}
We use this map for polynomials in the completed gap primitives and
their derivatives.  Theorem~\ref{gm:closure} specifies its continuation.
It also has an explicit ordinary-series meaning.  With
\(F(y,z)=y^{-s}z^{-t}P(y,z)\),
\begin{align}\label{gm:ordinary-E}
 \mathcal E_{a,b}^{s,t}[P]
 ={}&\sum_{n>j\ge0}F(n+b,j+b)
     +\sum_{0\le n\le j}F(n+a,j+a)\notag\\
    &-\sum_{n,j\ge0}F(n+a,j+b),
\end{align}
whenever these sums converge absolutely.  This follows by expanding
the depth-two case of \eqref{gm:relative}; the two \(a,a\) terms
combine to the weak opposite ordering.

Put
\begin{equation}\label{gm:gap}
 g_u(y,z)=Q(u,z+1)-Q(u,y).
\end{equation}
Taking the difference between the insertions at the two boundary
letters in Lemma~\ref{gm:one-tail}, and using \eqref{gm:Q}, gives
\begin{equation}\label{gm:local-insertion}
 \mathcal E_{a,b}^{s,t}[g_u]=\mathcal H_{a,b}(s,u,t).
\end{equation}
The same assertion holds with arbitrary additional boundary
decorations.

\begin{lemma}[Products of gap primitives]
\label{gm:ordered-partitions}
For independent complex \(u_1,\ldots,u_k\), let
\(U_B=\sum_{i\in B}u_i\).  Then
\begin{equation}\label{gm:ordered-partition-formula}
 \mathcal E_{a,b}^{s,t}\!\left[\prod_{i=1}^kg_{u_i}\right]
 =\sum_{(B_1,\ldots,B_\ell)}
      \mathcal H_{a,b}(s,U_{B_1},\ldots,U_{B_\ell},t),
\end{equation}
where the sum ranges over all ordered set partitions of
\(\{1,\ldots,k\}\).  The identity is jointly entire.
\end{lemma}

\begin{proof}
After some marked positions have been inserted, write their variables
as \(x_1,\ldots,x_\ell\), with \(x_0=y\) and
\(x_{\ell+1}=z\).  The exact telescoping identity
\begin{equation}\label{gm:split}
 g_v(y,z)=\sum_{h=0}^\ell g_v(x_h,x_{h+1})
                              +\sum_{h=1}^\ell x_h^{-v}
\end{equation}
follows from \eqref{gm:Q}.  It is an identity of primitive functions,
without an integer-lattice assumption on the external endpoints.
By \eqref{gm:local-insertion}, the first sum inserts the next label
in each possible strict slot, while the second merges it with each
existing block.  Starting with one marker, this constructs every
ordered set partition once.  All operations are identities between the
initial convergent quotients; Theorem~\ref{gm:closure} supplies their
joint continuation.
\end{proof}

For a set partition \(\pi\) of \(\{1,\ldots,k\}\), set
\[
 \mu(\pi)=\prod_{B\in\pi}(-1)^{|B|-1}(|B|-1)!.
\]

\begin{theorem}[Local symmetric-sum identity]
\label{gm:symmetric}
For arbitrary complex spectral arguments,
\begin{align}\label{gm:symmetric-formula}
 &\sum_{\sigma\in S_k}
       \mathcal H_{a,b}(s,u_{\sigma(1)},\ldots,u_{\sigma(k)},t)
 \notag\\
 &\qquad=\mathcal E_{a,b}^{s,t}\!\left[
        \sum_{\pi}\mu(\pi)\prod_{B\in\pi}g_{U_B}(y,z)\right].
\end{align}
\end{theorem}

\begin{proof}
Apply M\"obius inversion on the lattice of equality partitions to
\eqref{gm:ordered-partition-formula}.  A block \(B\) has M\"obius
coefficient \((-1)^{|B|-1}(|B|-1)!\).  This is the classical
symmetric-sum argument \cite{gm:Hoffman1992,gm:Hoffman2018};
Lemma~\ref{gm:ordered-partitions} proves that it acts on the present
local analytic expressions.  For example, at \(k=2\),
\[
 \mathcal H(s,u,v,t)+\mathcal H(s,v,u,t)
          =\mathcal E[g_ug_v-g_{u+v}].
\]
Every function used in the inversion has already been continued
jointly entirely.
\end{proof}

\subsection{Independent derivatives and identical marked clusters}

For integers \(r_i\ge0\) and centers \(\lambda_i\in\mathbb C\),
write
\[
 r_B=\sum_{i\in B}r_i,\qquad
 \lambda_B=\sum_{i\in B}\lambda_i,\qquad
 p_B(y,z)=Q^{[r_B]}(\lambda_B,z+1)
                              -Q^{[r_B]}(\lambda_B,y).
\]

\begin{corollary}[All independently marked symmetric sums]
\label{gm:marked}
The exact identity
\begin{align}\label{gm:marked-formula}
 &\left.\prod_{i=1}^k\partial_{u_i}^{r_i}
    \sum_{\sigma\in S_k}
      \mathcal H_{a,b}(s,u_{\sigma(1)},\ldots,u_{\sigma(k)},t)
                       \right|_{u_i=\lambda_i}\notag\\
 &\qquad=\mathcal E_{a,b}^{s,t}\!\left[
               \sum_\pi\mu(\pi)\prod_{B\in\pi}p_B(y,z)\right]
\end{align}
uses at most \(2^k-1\) different gap primitives.  At every merged
resonance \(\lambda_B=1\), the corresponding primitive is
\begin{equation}\label{gm:merged-stieltjes}
 p_B(y,z)=(-1)^{r_B}
                  \bigl(\gamma_{r_B}(z+1)-\gamma_{r_B}(y)\bigr).
\end{equation}
\end{corollary}

\begin{proof}
Differentiate the entire identity \eqref{gm:symmetric-formula}.
Within a block, every variable enters \(U_B\) with coefficient one,
so its derivative order is exactly \(r_B\), without an extra
multinomial coefficient.  Equation~\eqref{gm:stieltjes} proves
\eqref{gm:merged-stieltjes}.  Simultaneous resonances therefore require
no directional limit or separately chosen finite part.
\end{proof}

When all pairs \((\lambda_i,r_i)\) equal \((\lambda,r)\), each
permutation term is identical.  Put
\begin{align}\label{gm:Bell}
 p_j&=Q^{[jr]}(j\lambda,z+1)-Q^{[jr]}(j\lambda,y),\notag\\
 e_k(p_1,\ldots,p_k)
   &=[v^k]\exp\left(\sum_{j=1}^k\frac{(-1)^{j-1}}j p_jv^j\right).
\end{align}
The set-partition formula becomes the individual ordered identity
\begin{equation}\label{gm:identical}
 \left.\partial_{u_1}^r\cdots\partial_{u_k}^r
       \mathcal H_{a,b}(s,u_1,\ldots,u_k,t)
                         \right|_{u_1=\cdots=u_k=\lambda}
       =\mathcal E_{a,b}^{s,t}[e_k(p_1,\ldots,p_k)].
\end{equation}
These are independent derivatives taken before specialization.  A
total derivative in one common variable would give a different sum.
The coefficient in \eqref{gm:Bell} is a finite polynomial operation;
no convergence of an all-\(k\) marked generating series is asserted.
The first useful polynomials are
\begin{align}\label{gm:Bell-examples}
 e_2&=\tfrac12(p_1^2-p_2),\notag\\
 e_3&=\tfrac16(p_1^3-3p_1p_2+2p_3),\notag\\
 e_4&=\tfrac1{24}
       (p_1^4-6p_1^2p_2+3p_2^2+8p_1p_3-6p_4).
\end{align}

\subsection{Explicit Gamma and Stieltjes identities}

For zero-centered first derivatives define
\[
 L(y,z)=\log\Gamma(z+1)-\log\Gamma(y),\qquad
 R_j(y,z)=\zeta^{[j]}(0,z+1)-\zeta^{[j]}(0,y).
\]
The classical identity
\(\zeta^{[1]}(0,x)=\log\Gamma(x)-\tfrac12\log(2\pi)\)
\cite{DLMF} makes \(p_1=L\), while \(p_j=R_j\) for
\(j\ge2\).  The source's first two-marker target is therefore
\begin{equation}\label{gm:gamma-square}
 \boxed{\left.\partial_u\partial_v
       \mathcal H_{a,b}(s,u,v,t)\right|_{u=v=0}
       =\frac12\mathcal E_{a,b}^{s,t}[L^2-R_2].}
\end{equation}
This is an identity at arbitrary independent \(s,t\) and arbitrary
right-half-plane endpoints.  It supplies ordinary absolutely convergent
sums through \eqref{gm:ordinary-E} whenever \(\Re s,\Re t>3\):
\(\log\Gamma(x)=O(x\log x)\) and
\(\zeta^{[2]}(0,x)=O(x\log^2x)\) on each fixed lattice give
summability of every separated letter.  Elsewhere it is the specified
entire continuation.  The third and fourth identical marked clusters
are obtained from \eqref{gm:Bell-examples}; no new ordered coordinate
is required for these symmetric identical letters.

To connect this directly with normalized antiderivatives, write the
Hurwitz parameter-derivative law in its entire form
\[
 \partial_x\zeta(u,x)=-1-uQ(u+1,x).
\]
Differentiating at \(u=0\), for every \(j\ge1\), gives
\begin{equation}\label{gm:primitive}
 \partial_x\zeta^{[j]}(0,x)=(-1)^j j\gamma_{j-1}(x),\qquad
 R_j(y,z)=(-1)^j j\int_y^{z+1}\gamma_{j-1}(x)\,dx.
\end{equation}
The paths lie in the right half-plane; the constants are fixed by
vanishing when \(z+1=y\).  Thus \eqref{gm:gamma-square} is also
\begin{equation}\label{gm:gamma-integral}
 \left.\partial_u\partial_v\mathcal H_{a,b}(s,u,v,t)\right|_{0,0}
  =\frac12\mathcal E_{a,b}^{s,t}\!\left[
       L(y,z)^2-2\int_y^{z+1}\gamma_1(x)\,dx\right].
\end{equation}
The logarithm in \(L\) remains the difference of analytic log-Gamma
branches; it is not replaced by an independently chosen logarithm of
a ratio.

An example in which the merging itself creates a Stieltjes coefficient
uses two half-order marks.  Set
\(A(y,z)=\zeta^{[1]}(1/2,z+1)-\zeta^{[1]}(1/2,y)\).
Then
\begin{equation}\label{gm:half-pair}
 \left.\partial_u\partial_v\mathcal H_{a,b}(s,u,v,t)
                                      \right|_{u=v=1/2}
 =\frac12\mathcal E_{a,b}^{s,t}
          [A(y,z)^2-\gamma_2(z+1)+\gamma_2(y)].
\end{equation}
Indeed the merged index is one and its total derivative order is two.

For an unequal pair, the symmetric distinction is essential.  Define
\[
 J_{0,1}(s,t)=\left.\partial_u\partial_v
 \bigl\{\mathcal H_{a,b}(s,u,v,t)+\mathcal H_{a,b}(s,v,u,t)\bigr\}
                                                    \right|_{u=0,v=1}.
\]
The exact mixed Gamma--Stieltjes formula is
\begin{equation}\label{gm:mixed-pair}
 J_{0,1}(s,t)=\mathcal E_{a,b}^{s,t}\!\left[
 -L(y,z)\{\gamma_1(z+1)-\gamma_1(y)\}
                         -\gamma_2(z+1)+\gamma_2(y)\right].
\end{equation}
This determines the sum of the two ordered derivatives, not each
summand separately.

\subsection{Several gaps and the exact scope of the reduction}

For surviving variables \(x_1,\ldots,x_d\), the three types of
completed gaps are
\begin{equation}\label{gm:all-gaps}
 \begin{cases}
 Q(u,x_1+1)-Q(u,a),&\text{above the first variable},\\
 Q(u,x_{j+1}+1)-Q(u,x_j),&\text{between successive variables},\\
 Q(u,b)-Q(u,x_d),&\text{below the last variable}.
 \end{cases}
\end{equation}
On an absolute-convergence domain, $Z_c$ is a quasi-shuffle character
on lattice functions with pointwise multiplication. Characters with
values in a commutative algebra form a group under deconcatenation
convolution; therefore $Z_a^{-1}*Z_b$ has this property for decorated
letters as well. The continuation proved above preserves these
identities. The upper boundary in \eqref{gm:all-gaps} follows from the
one-tail lemma at the first position and the shift law. The lower
boundary follows by subtracting that lemma at the last position from
the singleton stuffle identity, using
$\mathcal H_{a,b}(u)=Q(u,b)-Q(u,a)$.

Apply the partition polynomial independently in each gap and multiply
the resulting polynomials before expanding into decorated letters.
The proof by insertion applies gap by gap, and
Theorem~\ref{gm:closure} gives joint continuation, including all shared
spectral resonances.  It determines the ordered word whenever each gap
contains one mark or an identical marked cluster.  With no surviving
variables, the sole gap is \(Q(u,b)-Q(u,a)\), and the construction
reduces to the classical symmetric-sum formula for the relative
interpolant.

The accompanying program verifies the local symmetric identity through
five independent labels by exact integer arithmetic on formal decorated
words.  All endpoint factors and shorter words cancel, leaving precisely
the \(k!\) permutations.  A separate numerical check uses the original
nested \(Z_c\)-series and relative quotient at nonintegral complex
endpoints: increasing the cutoff from 80 to 160 terms reduces its
absolute residual from \(1.66\times10^{-36}\) to
\(2.69\times10^{-40}\).  An independent reconstruction of the actual
zero-marker derivative at those noninteger endpoints, using convergent
nested sums and a symmetrized top-tail continuation, gives residuals
\(1.40\times10^{-27}\) and \(3.73\times10^{-31}\) at the same
cutoffs.  Individual singular double-zeta terms are never differentiated
at \((0,0)\).  Five direct finite-interval logarithmic
comparisons test negative centers, the Gamma square, three marks, and
merged Stieltjes resonances at 60 decimal digits; each residual is below
\(8\times10^{-62}\).  These are diagnostics of the algebra and
normalization, not interval certificates or substitutes for the
analytic proof.

The audited source correctly withheld the several-marker claim; no
correction to its single-marker theorem is needed.  The present theorem
answers its first Q4 target and the fully symmetric part of Q5.  A
remaining precise question is to represent the antisymmetric ordered
two-marker component with unequal marks.  Another is to classify which
unmarked nonpositive deletions preserve this symmetric span: such
deletions can shift different marked centers unequally, so their full
ordered extension is not implied by the Bell formula.
